\documentclass[a4paper,
	11pt,
	%draft,  %% deactivate graphics, mark long lines and use rulers in headings and footer
	%overfullrule=true,  % mark long lines
	]{scrartcl}

\usepackage[utf8]{inputenc}
\usepackage[ngerman]{babel}
\usepackage[draft=false]{scrlayer-scrpage} %the new srcpage2; part of Koma-Script; used for header and footer
		% draft=false  disables the rulers in headings and footer when draft is enabled in documentclass
\usepackage{courier}
\usepackage{mathpazo}
\usepackage[scaled=.95]{helvet}
\usepackage{graphicx}
\usepackage{fancyvrb,newverbs}
\usepackage[table]{xcolor} % load the table option to use \cellcolor and \rowcolors
\usepackage[colorlinks=true,
        linkcolor=black,
        citecolor=black,
        filecolor=black,
        urlcolor=blue,
        bookmarks=true,
        bookmarksopen=true,
        bookmarksopenlevel=3,
        plainpages=false,
        pdfpagelabels=true]{hyperref}
\usepackage{longtable,makecell}
\usepackage{tabularx}
\usepackage{dcolumn} % for aligning decimal points in tabulars
\usepackage{booktabs} % for better look and feel of regular tabular
\usepackage[%
%	format=plain,
%	%labelsep=newline,
%	singlelinecheck=false,
%	tableposition=top,  %this tells the captionsetup where skip between caption and float is inserted
%	figureposition=bottom,
%	skip=1cm,
	]{caption}  % for \captionsetup; customizable captions for floating environments
		    %  Loading caption to link to beginning of floats, not the caption.
%\captionsetup[table]{singlelinecheck=false,skip=1em} %%% But komascript already offers ways to customise the caption.
%\usepackage{color}
\usepackage{xspace}
\usepackage{microtype}
\usepackage{versions}
%\usepackage[nomarkers,figuresonly,noheads,nolists]{endfloat}
\usepackage[official]{eurosym}
		% use \texteuro for euro sign matching the current font.
		% use \euro for official euro sign
%\usepackage{marvosym} %prints official EUR symbol with \EUR
\usepackage[
	toc,
	acronym,
	%section=subsection, % print the glossary as a subsection
	section=section,     % print the glossary as a section
	numberedsection=autolabel,
	nopostdot=true]{glossaries} % after hyperref
\makeglossaries

\addto\captionsngerman{%  we use babel; so we can’t just redefine \acronymname. It has to be added to the captionslist for ngerman.
\renewcommand{\acronymname}{Abkürzungen}
}



\pagestyle{scrheadings} 
\setlength{\headheight}{2cm}
\setcounter{page}{0} 

% Durch Freischalten der folgenden Zeile wird die Nummerierung der Kapitel deaktiviert:
% \setcounter{secnumdepth}{-1}

\setkomafont{pagefoot}{\normalfont\scriptsize\sffamily}
\setkomafont{pagehead}{\rmfamily}
\setkomafont{subject}{\huge\sffamily}


\definecolor{b1blau}{RGB}{34,71,121}
\definecolor{b1orange}{RGB}{252,189,43}
\definecolor{greylist}{RGB}{224,224,224}

% Additional Colours from GIMP Palette
% Name: b1_farbpalette_210907

\definecolor{blue1}  {RGB}{42,79,131}
\definecolor{blue2}  {RGB}{63,102,155}
\definecolor{blue3}  {RGB}{81,122,176}
\definecolor{blue4}  {RGB}{87,129,188}
\definecolor{blue5}  {RGB}{139,180,234}
\definecolor{green1} {RGB}{73,157,28}
\definecolor{green2} {RGB}{122,192,66}
\definecolor{green3} {RGB}{147,225,66}
\definecolor{red1}   {RGB}{165,29,45}
\definecolor{red2}   {RGB}{224,27,36}
\definecolor{red3}   {RGB}{246,97,81}
\definecolor{orange1}{RGB}{250,140,31}
\definecolor{orange2}{RGB}{247,155,31}
\definecolor{grey1}  {RGB}{80,80,80}
\definecolor{b1grau} {RGB}{189,189,189}
\definecolor{grey2}  {RGB}{230,230,230}


\usepackage[final]{listings}
\lstset{
	frame=leftline,
	numbers=left,
	numberstyle=\tiny,
	stepnumber=1,
	breaklines=true,
	postbreak=\raisebox{0ex}[0ex][0ex]{\ensuremath{\hookrightarrow\space}},
	basicstyle=\small\ttfamily,
	showstringspaces=false,
  %~ extendedchars=true,
  %~ inputencoding=utf8,
  %~ frame=single,
  %~ framerule=0pt,
  %~ framesep=7pt,
  %~ backgroundcolor=\color{greylist},
  %~ literate={Ö}{{\"O}}1 {Ä}{{\"A}}1 {Ü}{{\"U}}1 {ß}{{\ss}}2 {ü}{{\"u}}1 {ä}{{\"a}}1 {ö}{{\"o}}1 {µ}{\textmu}1
}
\renewcommand{\lstlistlistingname}{Code- und Konfigurationsbeispiele} %% Erscheint im Inhaltsverzeichnis
\renewcommand{\lstlistingname}{Code/Konfiguration} %% Als Titel von Codebeispielen
\lstloadlanguages{%
	bash,
	Awk,
	HTML,
	Java,
	[5.3]Lua,
	make,
	go,%
	Python,
	PHP,
	Perl,
	Ruby,
	SQL,
	TeX,
	XML,
	}


\newenvironment{lcverbatim}
 {\SaveVerbatim{cverb}}
 {\endSaveVerbatim
  \flushleft\fboxrule=0pt\fboxsep=.5em
  \colorbox{greylist}{%
    \makebox[\dimexpr\linewidth-0\fboxsep][l]{\BUseVerbatim{cverb}}%
  }
  \endflushleft
}

%% B1-themed Hinweisboxen
\usepackage[tikz]{bclogo}
% Definitionen für Hinweisgrafiken
\renewcommand{\bcdanger}{\includegraphics[width=0.85\logowidth]{images/bc_danger_icon}}
\renewcommand{\bcattention}{\includegraphics[width=0.85\logowidth]{images/bc_attention_icon}}
\renewcommand{\bclampe}{\includegraphics[width=0.85\logowidth]{images/bc_lampe_icon}}
\renewcommand{\bcinfo}{\includegraphics[width=0.85\logowidth]{images/bc_info_icon}}


\newcommand{\todo}[1]{\textcolor{red}{#1}}
\newcommand{\beins}{B1 Systems GmbH}


% Definitionen für Hinweisboxen
%\newenvironment{name}[num]{before}{after}
% \def\boxesgerman{%
%  \def\warning{Warnung}%
%  \def\important{Wichtig}%
%  \def\tip{Tipp}%
%  \def\note{Hinweis}%
% }
% \def\boxesenglish{%
%  \def\warning{Warning}%
%  \def\important{Important}%
%  \def\tip{Tip}%
%  \def\note{Note}%
% }
% \ifthenelse{\equal{\configLang}{english}}{\boxesenglish}{\boxesgerman}

\newenvironment{b1warning}[1]{%
\begin{bclogo}[logo=\bcdanger,couleur=blue!10,barre=none,epBord=0.4,ombre=true,couleurOmbre=black!60,blur,arrondi=0.1,sousTitre=#1]{Achtung}}{\end{bclogo}}

\newenvironment{b1important}[1]{%
\begin{bclogo}[logo=\bcattention,couleur=blue!10,barre=none,epBord=0.4,ombre=true,couleurOmbre=black!60,blur,arrondi=0.1,sousTitre=#1]{Wichtig}}{\end{bclogo}}

\newenvironment{b1tip}[1]{%
\begin{bclogo}[logo=\bclampe,couleur=blue!10,barre=none,epBord=0.4,ombre=true,couleurOmbre=black!60,blur,arrondi=0.1,sousTitre=#1]{Tipp}}{\end{bclogo}}

\newenvironment{b1note}[1]{%
\begin{bclogo}[logo=\bcinfo,couleur=blue!10,barre=none,epBord=0.4,ombre=true,couleurOmbre=black!60,blur,arrondi=0.1,sousTitre=#1]{Hinweis}}{\end{bclogo}}

%%%%%%%%%%%%%%%%%%%%%%%%%%%%%%%%%%%%%%%%%%%%%%%%%%%%%%%%%%%%%%%
% Fußzeile                                                    %
%%%%%%%%%%%%%%%%%%%%%%%%%%%%%%%%%%%%%%%%%%%%%%%%%%%%%%%%%%%%%%%

% \cfoot{\scriptsize
% \hrule
% \vspace{0.3cm}
% \begin{footnotesize}\textbf{\textcolor{b1blau}{B}\textcolor{b1orange}{1} Systems GmbH \textemdash\xspace Linux/Open Source Consulting, Training, Managed Service \textcolor{b1orange}{\&} Support}\end{footnotesize}\\
% \vspace{0.07cm}
% \begin{footnotesize}Osterfeldstraße 7 / D-85088 Vohburg / +49 (0)8457 - 931096\end{footnotesize}\\
% \begin{footnotesize}info@\textcolor{b1blau}{b}\textcolor{b1orange}{1}-systems.de / www.\textcolor{b1blau}{b}\textcolor{b1orange}{1}-systems.de \end{footnotesize}\\
% }

\hypersetup{
%hidelinks,
%linkcolor=b1blau %% use this later, after toc
}
\cfoot{\scriptsize
\hrule
\vspace{0.3cm}
\begin{footnotesize}\textbf{\textcolor{b1blau}{B}\textcolor{b1orange}{1} Systems GmbH \textemdash\xspace Linux/Open Source Consulting, Training, Managed Service \textcolor{b1orange}{\&} Support}\end{footnotesize}\\
\vspace{0.07cm}
\begin{footnotesize}Osterfeldstraße 7 / D-85088 Vohburg / +49 (0)8457 - 931096\end{footnotesize}\\
\begin{footnotesize}
\href{mailto:info@b1-systems.de}{info@b1-systems.de} / \href{https://www.b1-systems.de}{www.b1-systems.de} \end{footnotesize}\\
}




%\setcounter{tocdepth}{2} % How many Depths schould be in the toc?



%%%%%%%%%%%%%%%%%%%%%%%%%%%%%%%%%
% LETS GEHT THIS PARTEY STARTET %
%%%%%%%%%%%%%%%%%%%%%%%%%%%%%%%%%

%%%%%%%%%%%%%%%%%%%%%%%%%%%%%%%%%%%%%%%%%%%%%%%%%%%%%%%%%%%%%%%
% Kopfzeile auf jeder Seite - BITTE ANPASSEN!                 %
%%%%%%%%%%%%%%%%%%%%%%%%%%%%%%%%%%%%%%%%%%%%%%%%%%%%%%%%%%%%%%%
\makeatletter
\ohead{%\@title
\small
	Die Erstellung eines Backends für eine Fahrzeugverwaltungssoftware\newline
\makeatother
Seite \pagemark}
\ihead{\includegraphics[height=1.5cm]{b1-systems.eps}}



\begin{document}
% Define glossary entries.
% Glossary (Glossar) are longer explanations of technical terms. Also plural forms are defined here.
%
% use the command \gls{} to print the entry into the text.
% use \glspl{} to print the plural form (if no plural is defined, just an 's' is added}
% use \Gls{} to print the singular form; but with a Capital Letter
% use \Glspl{} to print the plural form; but with a Capital Letter
% The last two are used for non-nouns appearing at the start of the sentence

\newglossaryentry{open source}  % define the entry
{
  name={Open Source}, % Appears in glossary
  text={Open Source}, % Appears in text
  plural={Open Source}, % Who would say such a thing?
  description=
    {
        Software ist Software, deren Quellcode öffentlich zugänglich ist, sodass Nutzer ihn einsehen, ändern und verteilen können.
    }
}

\newglossaryentry{rest api}  % define the entry
{
  name={REST API}, % Appears in glossary
  text={REST API}, % Appears in text
  plural={REST API}, % Who would say such a thing?
  description=
    {
        Eine REST-API (Representational State of Resource) ist eine Programmierschnittstelle, die es ermöglicht, Daten zwischen verschiedenen Systemen auszutauschen, indem sie Ressourcen über HTTP-Methoden wie GET, POST, PUT und DELETE anfordert und manipuliert.
    }
}

\newglossaryentry{linux}  % define the entry
{
  name={Linux}, % Appears in glossary
  text={Linux}, % Appears in text
  plural={Linux}, % Who would say such a thing?
  description=
    {
        ist ein freies und offenes Betriebssystem, das auf dem Linux-Kernel basiert und häufig in Servern, Desktops und eingebetteten Systemen verwendet wird.
    }
}

\newglossaryentry{vscode}  % define the entry
{
  name={Visual Studio Code}, % Appears in glossary
  text={Visual Studio Code}, % Appears in text
  plural={Visual Studio Code}, % Who would say such a thing?
  description=
    {
        ist ein leistungsstarker, kostenloser Quelltext-Editor von Microsoft, der eine Vielzahl von Programmiersprachen unterstützt und durch Erweiterungen anpassbar ist.
    }
}

\newglossaryentry{python}  % define the entry
{
  name={Python}, % Appears in glossary
  text={Python}, % Appears in text
  plural={Python}, % Who would say such a thing?
  description=
    {
        ist eine höhere Programmiersprache, die für ihre Einfachheit, Lesbarkeit und Vielseitigkeit bekannt ist und häufig für Webentwicklung, Datenanalyse, künstliche Intelligenz und Automatisierung verwendet wird.
    }
}

\newglossaryentry{fastapi}  % define the entry
{
  name={FastAPI}, % Appears in glossary
  text={FastAPI}, % Appears in text
  plural={FastAPI}, % Who would say such a thing?
  description=
    {
        ist ein modernes, leistungsstarkes und leichtgewichtiges Python-Framework für die Entwicklung von APIs (Application Programming Interfaces) mit einer starken Betonung auf Geschwindigkeit, Sicherheit und Benutzerfreundlichkeit.
    }
}

\newglossaryentry{pydantic}  % define the entry
{
  name={Pydantic}, % Appears in glossary
  text={Pydantic}, % Appears in text
  plural={Pydantic}, % Who would say such a thing?
  description=
    {
        ist eine Python-Bibliothek, die es ermöglicht, Datenmodelle zu definieren und zu validieren, indem sie eine einfache und intuitive Schnittstelle für die Erstellung von robusten, typsicheren und selbst dokumentierenden Datenstrukturen bietet.
    }
}

\newglossaryentry{crud}  % define the entry
{
  name={CRUD}, % Appears in glossary
  text={CRUD}, % Appears in text
  plural={CRUD}, % Who would say such a thing?
  description=
    {
        steht für Create, Read, Update und Delete und beschreibt die grundlegenden Operationen, die auf Daten in einer Datenbank oder einem Datenspeicher durchgeführt werden können.
    }
}

\newglossaryentry{fuzzy}  % define the entry
{
  name={Fuzzy}, % Appears in glossary
  text={Fuzzy}, % Appears in text
  plural={Fuzzy}, % Who would say such a thing?
  description=
    {
      Die Fuzzy-Suche ist eine Suchmethode, die ähnliche, aber nicht exakt übereinstimmende Begriffe findet, indem sie Tippfehler, unterschiedliche Schreibweisen oder ähnliche Wörter berücksichtigt, um die Suchergebnisse zu erweitern und relevanter zu gestalten.
    }
}

\newglossaryentry{gitlab}  % define the entry
{
  name={GitLab}, % Appears in glossary
  text={GitLab}, % Appears in text
  plural={GitLab}, % Who would say such a thing?
  description=
    {
      ist eine webbasierte Plattform zur Versionskontrolle und Zusammenarbeit an Softwareprojekten, die Funktionen wie Repository-Management, Continuous Integration und Projektmanagement in einer einzigen Anwendung vereint.
    }
}

% Define acronyms.
% Acronyms are just short expansions of shortened terms.
% They are only listed at the end if they occur in the document
% The first apperance in the text is printed with the long description and the short version in (parentheses)
%
% Acronyms (Or Abbreviations) have a "First-Use-Flag" which is handled by LaTeX internally
% When an entry that is defined with \newacronym is used with a \gls-like command this
% entry is put into the text with "Long Text (short)" and the First-Use-Flag is unset.
% All subsequent calls for this entry only insert the short version.
% To specifically use the short version without changing the Flag, use \glsentryshort{label}
% To unset the Flag use \glsunset{label}. To reset the Flag use \glsreset{label}.

%\newacronym{name-of-the-entry}{How it appears in the text}{What it means}

\newacronym{api}{API}{Application Programming Interface}
\newacronym{cli}{CLI}{Commandline Interface}
\newacronym{csv}{CSV}{Comma Seperated Values}
\newacronym{ui}{UI}{Benutzerschnittstelle (User Interface)}
\newacronym{url}{URL}{Uniform Resource Locator}
\newacronym{http}{HTTP}{Hyptertext Transfer Protocol}
\newacronym{https}{HTTPS}{Hyptertext Transfer Protocol Secure}
\newacronym{html}{HTML}{Hypertext Markup Language}
\newacronym{ssl}{SSL}{Secure Socket Layer}
\newacronym{tcp}{TCP}{Transmission Control Protocoll}
\newacronym{tls}{TLS}{Transport Layer Security}
\newacronym{xml}{XML}{erweiterbare Auszeichnungssprache (Extensible Markup Language)} % in this file all glossary entries and acronyms are defined


\hyphenation{
Netz-werk-na-mens-raum 
% pa-r-al-lel
SIG-HUP
Open-Source
Open-Source-Bereich
%Über-wach-ungs-ge-rä-te %Umlaut   %; muss manuell getrennt werden
An-wen-dungs-be-reich
}

% Keine Absatz-Einrückung, Abstand nach Absätzen
\parindent0pt
\parskip11pt

\begin{titlepage}
\titlehead{
\centering
\includegraphics[height=5cm]{b1-systems.eps}
}
%% FIXME: sobald RGB EPS da ist, bitte eps hinlegen!
%%%%%%%%%%%%%%%%%%%%%%%%%%%%%%%%%%%%%%%%%%%%%%%%%%%%%%%%%%%%%%%
% Bei \title den Titel für die Titelseite eintragen,          % 
% \author ist der Autor und \subject der Untertitel.          %
%%%%%%%%%%%%%%%%%%%%%%%%%%%%%%%%%%%%%%%%%%%%%%%%%%%%%%%%%%%%%%% 
\title{Die Erstellung eines Backends für eine Fahrzeugverwaltungssoftware bei der \beins{}}
\author{Moritz Stefan Menzel}
\publishers{\beins}
\subject{Prüflingsnummer: xxxxxx}
\maketitle\thispagestyle{empty}
\end{titlepage}

%%%%%%%%%%%%%%%%%%%%%%%%%%%%%%%%%%%%%%%%%%%%%%%%%%%%%%%%%%%%%%%
% Diese Tabelle enthält die Versionshistorie.                 % 
% Aufbau von Tabellen: Jedes Feld wird mit & abgeschlossen,   % 
% jede Zeile mit  \\                                          %
% ein \hline sorgt für eine horizontale Trenn-Linie.          %
%%%%%%%%%%%%%%%%%%%%%%%%%%%%%%%%%%%%%%%%%%%%%%%%%%%%%%%%%%%%%%% 
%~ \section*{Versionshistorie}
%~ \begin{tabularx}{\textwidth}{|l|l|X|l|}
%~ \hline
%~ \textsc{\textbf{Version}} & \textsc{\textbf{Datum}} & \textsc{\textbf{Geändert wurde}} & \textsc{\textbf{Von wem}} \\
%~ \hline
%~ 00.1 & yyyy-mm-dd & new file 										& Some Dude \\
%% More entries
%~ 00.5 & 2018-05-06 & Kapitel Datenbank und Parallel Betrieb			& Some Dude \\
%~ 00.6 & 2018-05-11 & Umstrukturierung, Entknotung, logischer Aufbau	& Some Dude \\
%~ 00.7 & 2018-05-14 & Lektorat											& janke, et.al. \\
%~ \hline
%~ \end{tabularx}

%~ \pagebreak

%%%%%%%%%%%%%%%%%%%%%%%%%%%%%%%%%%%%%%%%%%%%%%%%%%%%%%%%%%%%%%%
% Inhaltsverzeichnis                                          %
%%%%%%%%%%%%%%%%%%%%%%%%%%%%%%%%%%%%%%%%%%%%%%%%%%%%%%%%%%%%%%%
\pagenumbering{roman}  %switch to lowercase roman letters pagenumbers on the toc
\tableofcontents  % Prints the Table of Contents

\pagebreak

\listoffigures   % Prints a list of images used in the document

\listoftables    % Prints a list of tables used in the document

\lstlistoflistings  % Prints a list of code examples....

\cleardoubleoddpage  % make sure the next page is a right page and insert empty page if necessary
\pagenumbering{arabic}  % reset page counter and use arabic pagenumberung for main part

\hypersetup{
%hidelinks,
linkcolor=b1blau %% use this after toc
}

%%%%%%%%%%%%%%%%%%%%%%%%%%%%%%%%%%%%%%%%%%%%%%%%%%%%%%%%%%%%%%%
% Inhalt                                                      %
%%%%%%%%%%%%%%%%%%%%%%%%%%%%%%%%%%%%%%%%%%%%%%%%%%%%%%%%%%%%%%%

%% Either start here with your tex(t), or use external
%% .tex-files relative to kundendoku.tex.
%% NOTE: don’t use the file extension (.tex) with \input

%%%%%%%%%%%%%%%%%%%%%%%%%%%%%%%%%%%%%
% Einleitung und Beschreibung des Projekts, Überblick über das Kundensystem
%%%%%%%%%%%%%%%%%%%%%%%%%%%%%%%%%%%%%

\section{Einleitung}

Die vorliegende Dokumentation beschreibt den Ablauf des IHK-Abschlussprojekts, das von dem Auszubildenden Moritz Stefan Menzel im Rahmen seiner Ausbildung zum Fachinformatiker in der Fachrichtung Anwendungsentwicklung durchgeführt wird.

\subsection{Unternehmensvorstellung}

Die \beins{} ist ein mittelständisches Trainings- und Consultingunternehmen, das sich auf \gls{open source}-Lösungen spezialisiert. Auf internationaler Ebene unterstützt die \beins{} ihre Kunden bei der Planung, Implementierung und Betreuung von IT-Infrastrukturen. Darüber hinaus bietet das Unternehmen eine Vielzahl von Schulungen in verschiedenen Informationstechnologien an, darunter Linux-Grundlagen, Kubernetes, Keycloak und viele weitere Themen.

Das umfangreiche Technologieportfolio der \beins{} umfasst unter anderem die Bereiche Virtualisierung, Container, Hochverfügbarkeit, Monitoring und Konfigurationsmanagement. Zudem verfügt das Unternehmen über eine Entwicklungsabteilung, die sowohl an \gls{open source}-Projekten mitwirkt als auch interne und externe Projekte realisiert.
Stand November 2025 beschäftigt die \beins{} etwa 200 Mitarbeiter in 5 Ländern.

\subsection{Projektbeschreibung}

Die \beins{} stellt ihren Mitarbeitern Fahrzeuge zur Verfügung, da sie regelmäßig zu Kundenbesuchen unterwegs sind. Darüber hinaus bietet das Unternehmen auch ein internes Auto-Leasing-Programm an, um den individuellen Bedürfnissen ihrer Mitarbeiter gerecht zu werden.

Die Fahrzeugdaten werden bislang handschriftlich auf Papier festgehalten und in Ordnern archiviert. Im Rahmen der Digitalisierung des Unternehmens und der Implementierung effizienterer Prozesse wird beschlossen, eine Fahrzeugverwaltungssoftware zu entwickeln, die den Arbeitsaufwand minimiert und die bisherige Methode ablöst.

\subsection{Projektziel}

Die Aufgabe des Auszubildenden besteht darin, ein Backend für die Fahrzeugverwaltungssoftware zu entwickeln.

Die Implementierung einer \gls{rest api}-Anwendung ermöglicht die umfassende Verwaltung von Fahrzeugdaten in einer Datenbank. Diese Fahrzeugdaten umfassen spezifische Informationen wie Marke, Kennzeichen sowie Vor- und Nachname des Mitarbeiters, dem das Fahrzeug zugeordnet ist.

Die \gls{api} ist für die grundlegenden Operationen zuständig: Erstellen, Abrufen, Aktualisieren und Löschen dieser Daten. Das zugrunde liegende Datenmodell definiert die Struktur der Fahrzeugdaten und stellt sicher, dass alle relevanten Informationen konsistent und korrekt gespeichert werden.

In Kombination mit der \gls{api}-Anwendung können Benutzer über verschiedene \gls{api}-Endpunkte auf die Fahrzeugdaten zugreifen. Diese Endpunkte sind spezifische URLs, die es ermöglichen, die oben genannten Operationen durchzuführen, und bieten eine klare Schnittstelle zur Interaktion mit den Fahrzeugdaten.

Zusammengefasst ermöglicht die \gls{rest api}-Anwendung eine effiziente Verwaltung von Fahrzeugdaten, wobei das Datenmodell die Struktur dieser Daten festlegt und die \gls{api}-Endpunkte den Zugriff und die Manipulation der Daten erleichtern.

\subsection{Projektumfeld}

Die Entwicklung der Fahrzeugverwaltungssoftware ist ein internes Projekt, das von der Geschäftsleitung der \beins{} initiiert wird. Herr Bernd Ritter, Linux-Entwickler bei der \beins{}, übernimmt die technische Verantwortung. Im Rahmen der agilen Entwicklungsweise gibt er wertvolles Feedback und führt Codereviews durch, um technische Fragen zu klären.

Die Umsetzung des Projekts erfolgt in den Büros der \beins{} in Rockolding auf einem firmeninternen Laptop mit einem \gls{linux} Betriebssystem. Ich erstelle den Code in \gls{vscode}, das als Entwicklungsumgebung dient. Für die Versionsverwaltung nutze ich die firmeninterne Instanz der Dev-Ops-Plattform \gls{gitlab}.

\subsection{Projektabgrenzung}

Folgende Einschränkungen werden im Vorfeld getroffen:
\begin{itemize}
    \item Das Projekt umfasst ausschließlich die Entwicklung des Backends für die Fahrzeugverwaltungssoftware.
    \item Die Entwicklung erfolgt in der Programmiersprache \gls{python}, wobei das Framework \gls{fastapi} und die Bibliothek \gls{pydantic} verwendet werden.
    \item Insgesamt entsteht ein Kleinstes realisierbares Produkt als Endergebnis.
\end{itemize}

\section{Projektplanung}

\subsection{Ist-Zustand} \label{Ist-Zustand}

Die Fahrzeugdaten werden derzeit handschriftlich in einem Ordner archiviert, was zu mehreren Problemen führt:

\begin{itemize}
    \item Die Daten sind nicht digital erfasst, was eine einfache Durchsuchung und Auswertung unmöglich macht.
    \item Der Zugriff auf die Fahrzeugdaten ist ausschließlich vor Ort möglich, was die Flexibilität erheblich einschränkt.
    \item Die Lesbarkeit der Handschrift ist oft unzureichend, was zu Missverständnissen führt.
    \item Die manuelle Pflege der Daten erfordert viel Zeit und ist anfällig für Fehler.
    \item Die Daten sind nicht ausreichend geschützt, was zu Beschädigungen oder dem Verlust der Dokumente führt.
\end{itemize}

Im Rahmen der Wirtschaftlichkeitsrechnung analysiere ich die Kosten des bisherigen Verfahrens gründlich. Die identifizierten Probleme bei der handschriftlichen Archivierung der Fahrzeugdaten verdeutlichen die Notwendigkeit, effizientere und digitalisierte Prozesse zu implementieren, um Zeit und Ressourcen zu sparen.

\subsection{Soll-Konzept} \label{Soll-Konzept}

In Abstimmung mit dem Auftraggeber werden folgende funktionale und nichtfunktionale Anforderungen erarbeitet:

Die Fahrzeugdaten werden über eine \gls{api} in einer Datenbank gespeichert, wobei die \gls{crud}-Methoden verwendet werden. Dies geschieht auf Grundlage eines vordefinierten Datenmodells.

Das Datenmodell legt fest, dass das Kennzeichen einzigartig ist. Zusätzlich erfasst es die Automarke sowie den Vornamen und Nachnamen des Mitarbeiters.

Die \gls{api} ermöglicht es, Fahrzeuge durch eine \gls{fuzzy}-Suche anhand des Kennzeichens, der Automarke, des Vornamens oder des Nachnamens zu finden. Darüber hinaus bietet sie Funktionen zum Bearbeiten, Löschen und Hinzufügen neuer Fahrzeuge.

Das Ergebnis ist in Form eines Aktivitätsdiagramms auf Seite~\pageref{Aktivitätsdiagramm} im Anhang~\ref{Aktivitätsdiagramm} zu finden.

\subsection{Zeitplanung}

Für dieses Projekt wählen wir die hybride Agile-Wasserfall-Softwareentwicklungsmethode, da einerseits die meisten Anforderungen zu Beginn bekannt sind und wir andererseits auf mögliche Änderungen seitens des Auftraggebers flexibel reagieren können. Zudem ermöglicht uns dies, in bestimmten Phasen zu iterieren. Durch diese Vorgehensweise können wir in ständigem Austausch mit dem Auftraggeber bleiben und wiederholt Feedback vom Projektbetreuer einholen, um in diesem Zeitraum die bestmögliche Lösung zu erarbeiten.

Die gewählte Methode zur Softwareentwicklung führt dazu, dass die Änderungen und der zeitliche Rahmen nicht chronologisch aufgelistet sind, wie in Tabelle~\ref{projekt-zeit} auf Seite~\pageref{projekt-zeit} dargestellt.

Im Projektantrag hat der Autor versehentlich 8 Stunden für die Erstellung des Soll-Konzepts angegeben. Diese Angabe wurde in der Zeittabelle korrigiert. Die tatsächliche Erstellung des Soll-Konzepts erforderte mehr Zeit, da eine sorgfältige Planung unerlässlich war. Auch die Implementierung des Backends nahm mehr Aufwand in Anspruch, da eine umfassende Einarbeitung in die verwendeten Frameworks notwendig war, um eine erfolgreiche Umsetzung zu gewährleisten. Darüber hinaus wurde weniger Zeit für Tests benötigt, da das hervorragend umgesetzte UML-Diagramm in der Erstellung des Soll-Konzepts erheblich zur Effizienz der Tests beigetragen hat.

\begin{table}[!h]
    \centering
    \caption{Erfasste Zeiten und Änderungen zum Projektantrag}
    \label{projekt-zeit}
    \begin{tabular}{lc<{\cellcolor{b1orange!50}}c<{\cellcolor{b1orange!70}}c}
    \toprule
        \cellcolor{white} & \multicolumn{2}{c}{\cellcolor{b1blau}\textcolor{white}{Zeit}} & \cellcolor{white} \\
        \cmidrule(lr){2-3}
        \scriptsize Tätigkeit & \scriptsize Geplant & \scriptsize Tatsächlich & \scriptsize Veränderung \\
        \midrule
        Ist-Analyse                           & 6 h  & 5 h                        & \cellcolor{green2!70}{$-$ 1 h} \\
        Erstellung Soll-Konzept               & 8 h  & 11 h                       & \cellcolor{red2!60}{$+$ 3 h} \\
        Wirtschaftlichkeitsrechnung           & 4 h  & 2 h                        & \cellcolor{green2!70}{$-$ 2 h} \\
        Implementierung Backend               & 29 h & 30 h                       & \cellcolor{red2!60}{$+$ 1 h} \\
        Testing und Debugging                 & 15 h & 13 h                       & \cellcolor{green2!70}{$-$ 2 h} \\
        Projektdokumentation                  & 18 h & 19 h                       & \cellcolor{red2!60}{$+$ 1 h} \\
        \midrule
            \bfseries{Gesamt}                 & \multicolumn{2}{c}{\cellcolor{b1blau}\textcolor{white}{80 h}} & \cellcolor{white}\\
    \bottomrule
    \end{tabular}
\end{table}

\subsection{Wirtschaftlichkeit}

\input{030-projektdurchfuehrung}

\input{050-fazit}

\section{Anhang} \label{Anhang}
\subsection{Aktivitätsdiagramm} \label{Aktivitätsdiagramm}

%%
%% \input{030-durchfuehrung}
%% 
%% \input{031-clientconnect}
%% \input{033-fluentd}
%% \input{035-reverseproxy}
%% \input{037-db}
%% \input{038-parallel}
%% 
%% \input{040-fazit}

\clearpage

\pagenumbering{Roman}  % use Roman upper case letters for glossary
\printglossaries % Prints your glossary (and by extent the acronyms). More configuration where the glossary package is loaded

%% If you want to insert the Acronyms and Glossary individually,
%% use these two instead of \printglossaries
% \setglossarysection{subsection}  % If one should be a subsection
%\printglossary[type=\acronymtype, style=listgroup]
%\printglossary[type=\glsdefaulttype, style=listgroup]



\end{document}
